\item A los átomos del mismo elemento pero con diferente número de neutrones en el núcleo se les denomina isótopos. El hidrógeno gaseoso ordinario es una mezcla de dos isótopos: hidrógeno-1 con núcleo protón, e hidrógeno-2 (deuterio) con núcleo deuterón (un protón y un neutrón enlazados).
\begin{enumerate}
    \item Use la ecuación dada en el problema 12 para demostrar que la diferencia en longitud de onda entre las líneas espectrales del hidrógeno-1 y del deuterio asociadas con una transición particular está dada por:
    $$ \Delta\lambda = \lambda \left( 1 - \frac{\mu_{\text{D}}}{\mu_{\text{H}}} \right) $$
    \item Evalúe la diferencia en longitud de onda para la línea $\text{H}_\alpha$ Balmer del hidrógeno, con longitud de onda $656.3\text{ nm}$. Harold Urey observó esta diferencia en 1931 y confirmó el descubrimiento del deuterio.
\end{enumerate}